\documentclass[12pt,a4paper]{article}
\usepackage[utf8]{inputenc}
\usepackage[T1]{fontenc}
\usepackage[indonesian]{babel}
\usepackage{geometry}
\usepackage{longtable}
\usepackage{array}
\usepackage{booktabs}
\usepackage{enumitem}
\usepackage[hidelinks,breaklinks]{hyperref}
\geometry{margin=2.5cm}

\begin{document}

\begin{center}
\textbf{\Large RENCANA PEMBELAJARAN SEMESTER (RPS)}\\
\textbf{Berbasis Outcome-Based Education (OBE)}\\[0.5cm]
\textbf{PROGRAM STUDI S1 INFORMATIKA}\\
\textbf{FAKULTAS TEKNIK}\\
\textbf{UNIVERSITAS LOREM IPSUM}\\
\end{center}

\vspace{1cm}

% ============================================================
% 1. IDENTITAS MATA KULIAH
% ============================================================
\section*{1. Identitas Mata Kuliah}
\begin{tabular}{ll}
Nama Program Studi & : S1 Informatika \\
Nama Mata Kuliah & : Pemrograman Game \\
Kode Mata Kuliah & : TIF-3XX \\
Semester & : 3 (Tiga) \\
SKS / Bobot Kredit & : 3 SKS (90\% Teori, 10\% Praktek) \\
Dosen Pengampu & : Dr. Lorem Ipsum, M.Kom. \\
\end{tabular}

\vspace{0.5cm}

% ============================================================
% 2. CAPAIAN PEMBELAJARAN LULUSAN (CPL)
% ============================================================
\section*{2. Capaian Pembelajaran Lulusan (CPL)}

CPL yang dibebankan pada mata kuliah ini dinyatakan dalam bentuk kompetensi lulusan (pengetahuan, keterampilan, sikap):

\begin{itemize}[leftmargin=*]
  \item \textbf{CPL-1 (Pengetahuan):} Menguasai konsep pengembangan game, arsitektur game engine, dan prinsip desain game secara mendalam.
  
  \item \textbf{CPL-2 (Keterampilan Umum):} Mampu menerapkan pemikiran logis, kritis, sistematis, dan inovatif dalam konteks pengembangan game dengan memperhatikan nilai humaniora.
  
  \item \textbf{CPL-3 (Keterampilan Khusus):} Mampu mengembangkan game 3D menggunakan Unity, menerapkan scripting C\#, dan mengintegrasikan komponen game (fisika, collider, UI, audio).
  
  \item \textbf{CPL-4 (Sikap):} Menunjukkan sikap bertanggung jawab atas pekerjaan di bidang keahliannya secara mandiri, mampu bekerja sama dalam tim, dan menerapkan etika dalam pengembangan game.
\end{itemize}

\vspace{0.5cm}

% ============================================================
% 3. CAPAIAN PEMBELAJARAN MATA KULIAH (CPMK)
% ============================================================
\section*{3. Capaian Pembelajaran Mata Kuliah (CPMK)}

Kemampuan atau kompetensi spesifik yang diharapkan mahasiswa kuasai setelah menyelesaikan mata kuliah:

\begin{itemize}[leftmargin=*]
  \item \textbf{CPMK-1:} Mahasiswa mampu memahami dan menjelaskan konsep pengembangan game serta arsitektur Unity sebagai game engine.
  
  \item \textbf{CPMK-2:} Mahasiswa mampu mengoperasikan Unity Editor dan mengatur Game Objects serta komponen (Transform, Rigidbody, Collider).
  
  \item \textbf{CPMK-3:} Mahasiswa mampu memprogram perilaku game dengan C\# script meliputi input player, fisika, dan collision detection.
  
  \item \textbf{CPMK-4:} Mahasiswa mampu membangun game 3D sederhana (Roll a Ball) dari awal hingga selesai dengan fitur lengkap.
  
  \item \textbf{CPMK-5:} Mahasiswa mampu membuat UI, mengelola game state, dan melakukan build/export game untuk platform target.
\end{itemize}

\vspace{0.5cm}

% ============================================================
% 4. SUB-CPMK / INDIKATOR PENCAPAIAN
% ============================================================
\section*{4. Sub-CPMK / Indikator Pencapaian}

Penjabaran CPMK menjadi indikator yang lebih terukur dan dapat diuji:

\begin{itemize}[leftmargin=*]
  \item \textbf{Sub-CPMK 1.1:} Menjelaskan peran game engine dalam pengembangan game
  \item \textbf{Sub-CPMK 2.1:} Menggunakan Inspector, Hierarchy, dan Scene View dalam Unity Editor
  \item \textbf{Sub-CPMK 2.2:} Mengatur Game Objects, Components, dan Transform
  \item \textbf{Sub-CPMK 3.1:} Menulis script MonoBehaviour dengan Update/FixedUpdate
  \item \textbf{Sub-CPMK 3.2:} Mengimplementasikan input player dan kontrol kamera
  \item \textbf{Sub-CPMK 3.3:} Menggunakan Rigidbody dan deteksi collision/trigger
  \item \textbf{Sub-CPMK 4.1:} Membuat collectibles dan sistem skor
  \item \textbf{Sub-CPMK 4.2:} Mengimplementasikan kondisi menang dan game over
  \item \textbf{Sub-CPMK 5.1:} Membuat UI Canvas dan menampilkan skor
  \item \textbf{Sub-CPMK 5.2:} Melakukan build game untuk platform target (Windows/standalone)
\end{itemize}

\vspace{0.5cm}

% ============================================================
% 5. MATERI PEMBELAJARAN (BAHAN KAJIAN)
% ============================================================
\section*{5. Materi Pembelajaran (Bahan Kajian)}

Daftar topik materi yang relevan dengan Sub-CPMK dan CPMK. Materi saling terhubung dalam satu proyek game ``Roll a Ball'' dari pertemuan ke pertemuan:

\begin{enumerate}[leftmargin=*]
  \item Pengenalan Pemrograman Game dan Game Engine
  \item Unity Editor dan Interface
  \item Game Objects, Components, dan Transform
  \item C\# Scripting Dasar untuk Unity (MonoBehaviour, variabel, input)
  \item Rigidbody dan Physics
  \item Input System dan Kontrol Player
  \item Camera Control dan Follow System
  \item Collision Detection dan Trigger
  \item Collectibles dan Game Logic
  \item UI Canvas dan Text
  \item Game State Management
  \item Prefabs dan Object Management
  \item Audio dan Visual Effects
  \item Build dan Deployment
\end{enumerate}

\textbf{Pemetaan 16 Pertemuan:}

\begin{longtable}{|c|>{\raggedright\arraybackslash}p{5.5cm}|c|>{\raggedright\arraybackslash}p{4cm}|}
\hline
\textbf{Pert.} & \textbf{Materi} & \textbf{Teori/Praktek} & \textbf{Keterkaitan Roll a Ball} \\
\hline
\endfirsthead
\hline
\textbf{Pert.} & \textbf{Materi} & \textbf{Teori/Praktek} & \textbf{Keterkaitan Roll a Ball} \\
\hline
\endhead
1 & Pengenalan Pemrograman Game, Unity, dan Game Engine & Teori & Orientasi proyek \\
2 & Unity Editor, Interface, dan Game Objects & Teori & Persiapan \\
3 & C\# Scripting Dasar untuk Unity & Teori & Fondasi scripting \\
4 & Setup Project ``Roll a Ball'' dan Play Area & Praktek (Unity Learn) & Pembuatan arena \\
5 & Moving the Player (Rigidbody, Input, Script) & Teori & Kontrol bola \\
6 & Moving the Camera (Follow target) & Teori & Kamera mengikuti \\
7 & Collision, Physics, dan Trigger & Teori & Landasan mekanik \\
8 & Creating Collectibles dan Deteksi Tabrakan & Teori & Item collectible \\
9 & \textbf{UTS} & -- & -- \\
10 & UI dan Canvas - Menampilkan Skor & Teori & Display skor \\
11 & Game State dan Win Condition & Teori & Kondisi menang \\
12 & Prefabs dan Object Management & Teori & Manajemen objek \\
13 & Audio dan Visual Polish & Teori & Efek suara dan visual \\
14 & Build dan Export Game & Praktek (Unity Learn) & Packaging game \\
15 & Refleksi dan Tanya Jawab & Teori & Presentasi/refleksi \\
16 & \textbf{UAS} & -- & -- \\
\hline
\end{longtable}

\textit{Catatan:} Pertemuan praktek (Pert. 4 dan 14) mengacu pada tutorial resmi Unity Learn Roll a Ball: \url{https://learn.unity.com/course/roll-a-ball}

\vspace{0.5cm}

% ============================================================
% 6. METODE PEMBELAJARAN
% ============================================================
\section*{6. Metode Pembelajaran}

Strategi atau pendekatan pembelajaran yang dipilih (mis. diskusi, PBL, proyek, praktik) sesuai OBE yang menekankan aktivitas mahasiswa. Komposisi pembelajaran: 90\% teori, 10\% praktek:

\begin{itemize}[leftmargin=*]
  \item \textbf{Ceramah dan Diskusi:} Penyampaian konsep pengembangan game, arsitektur Unity, dan prinsip desain game
  \item \textbf{Studi Kasus:} Analisis tutorial resmi Unity ``Roll a Ball'' sebagai referensi konsep
  \item \textbf{Self-Directed Practice:} Mahasiswa diarahkan ke Unity Learn untuk praktek membuat game ``Roll a Ball'' secara mandiri di luar jam kuliah. Sumber: \url{https://learn.unity.com/course/roll-a-ball}
  \item \textbf{Demonstrasi:} Dosen menunjukkan konsep dan contoh di Unity (terbatas, untuk mendukung pemahaman teori)
  \item \textbf{Tanya Jawab:} Diskusi konsep dan pemecahan masalah secara teori
\end{itemize}

\vspace{0.5cm}

% ============================================================
% 7. PENGALAMAN BELAJAR MAHASISWA
% ============================================================
\section*{7. Pengalaman Belajar Mahasiswa}

Deskripsi tugas, aktivitas, atau pengalaman belajar yang mendukung pencapaian Sub-CPMK:

\begin{itemize}[leftmargin=*]
  \item Membaca dan menganalisis materi teori pengembangan game serta arsitektur Unity
  \item Berpartisipasi dalam diskusi kelas dan tanya jawab konsep
  \item Mengerjakan kuis teori untuk evaluasi pemahaman
  \item Praktek: menyelesaikan tutorial game ``Roll a Ball'' di Unity Learn (\url{https://learn.unity.com/course/roll-a-ball}) sebagai tugas mandiri di luar jam kuliah
  \item Mempresentasikan atau merefleksikan hasil praktek (bila ada) di depan kelas
  \item Membuat dokumentasi singkat perkembangan praktek (jika menyelesaikan tutorial)
\end{itemize}

\vspace{0.5cm}

% ============================================================
% 8. KRITERIA, INDIKATOR, DAN BOBOT PENILAIAN
% ============================================================
\section*{8. Kriteria, Indikator, dan Bobot Penilaian}

Teknik/alat asesmen (tes, proyek, portofolio) dipetakan ke Sub-CPMK/CPMK. Bobot penilaian ditentukan sesuai kontribusinya terhadap CPL/CPMK. Bobot disesuaikan dengan komposisi pembelajaran 90\% teori dan 10\% praktek.

\begin{longtable}{|>{\raggedright\arraybackslash}p{2.5cm}|>{\raggedright\arraybackslash}p{4cm}|>{\raggedright\arraybackslash}p{5.5cm}|c|}
\hline
\textbf{Komponen} & \textbf{Teknik Asesmen} & \textbf{Indikator/CPMK} & \textbf{Bobot (\%)} \\
\hline
\endfirsthead
\hline
\textbf{Komponen} & \textbf{Teknik Asesmen} & \textbf{Indikator/CPMK} & \textbf{Bobot (\%)} \\
\hline
\endhead
\hline
\endfoot

Kuis Teori & Multiple Choice \& Essay & Sub-CPMK 1.1, 2.1, 3.1 & 20 \\
\hline
UTS & Written Exam \& Analisis & CPMK-1, CPMK-2, CPMK-3 & 25 \\
\hline
Tugas Praktek & Roll a Ball via Unity Learn + Checklist & Sub-CPMK 2.2, 3.2, 3.3, 4.1 & 10 \\
\hline
Proyek Game & Hasil praktek Unity Learn + Dokumentasi & CPMK-4, CPMK-5 & 5 \\
\hline
UAS & Comprehensive Exam & CPMK-1 s.d. CPMK-5 & 40 \\
\hline
\textbf{Total} & & & \textbf{100} \\
\hline
\end{longtable}

\textbf{Kriteria Penilaian:}
\begin{itemize}[leftmargin=*]
  \item A (85-100): Menguasai semua CPMK dengan sangat baik, mampu menerapkan dalam kasus kompleks
  \item B (70-84): Menguasai sebagian besar CPMK dengan baik
  \item C (60-69): Menguasai CPMK dasar dengan cukup
  \item D (50-59): Menguasai sebagian kecil CPMK
  \item E (<50): Belum menguasai CPMK yang ditetapkan
\end{itemize}

\vspace{0.5cm}

% ============================================================
% 9. EVALUASI DAN REFLEKSI PEMBELAJARAN (OPSIONAL)
% ============================================================
\section*{9. Evaluasi dan Refleksi Pembelajaran}

Penilaian sumatif/formatif untuk memantau ketercapaian outcome secara menyeluruh:

\begin{itemize}[leftmargin=*]
  \item \textbf{Evaluasi Formatif:} Checklist progress tiap pertemuan, kuis untuk memberikan feedback berkelanjutan
  \item \textbf{Evaluasi Sumatif:} UTS dan UAS untuk mengukur pencapaian CPMK secara komprehensif
  \item \textbf{Refleksi Mahasiswa:} Jurnal singkat perkembangan proyek game untuk refleksi diri
  \item \textbf{Evaluasi Dosen:} Survey kepuasan mahasiswa di tengah dan akhir semester
  \item \textbf{Continuous Improvement:} Analisis hasil penilaian untuk perbaikan RPS di semester berikutnya
\end{itemize}

\vspace{0.5cm}

% ============================================================
% 10. DAFTAR REFERENSI
% ============================================================
\section*{10. Daftar Referensi}

Sumber belajar utama yang digunakan dalam penyusunan materi dan asesmen:

\begin{enumerate}[leftmargin=*]
  \item Unity Technologies. (2024). \textit{Roll a Ball Tutorial}. Unity Learn. Retrieved from \url{https://learn.unity.com/course/roll-a-ball}
  
  \item Unity Technologies. (2024). \textit{Unity User Manual}. Retrieved from \url{https://docs.unity3d.com/Manual/}
  
  \item Unity Technologies. (2024). \textit{Unity Scripting API}. Retrieved from \url{https://docs.unity3d.com/ScriptReference/}
  
  \item Blackman, S. (2019). \textit{Beginning 3D Game Development with Unity 4} (2nd ed.). Apress.
  
  \item Hocking, J. (2018). \textit{Unity in Action: Multiplatform Game Development in C\#} (2nd ed.). Manning Publications.
  
  \item Goldstone, W. (2014). \textit{Unity Game Development in 24 Hours}. Sams Publishing.
  
  \item Microsoft. (2024). \textit{C\# Documentation}. Retrieved from \url{https://learn.microsoft.com/en-us/dotnet/csharp/}
  
  \item Unity Technologies. (2024). \textit{3D Beginner Game: Roll-a-Ball}. Unity Learn Project. Retrieved from \url{https://learn.unity.com/project/roll-a-ball}
\end{enumerate}

\vspace{1cm}

\begin{flushright}
\begin{tabular}{c}
Disusun oleh,\\[2cm]
\textbf{Dr. Lorem Ipsum, M.Kom.}\\
NIP. 123456789012345678\\
Tanggal Penyusunan: 9 Februari 2026
\end{tabular}
\end{flushright}

\end{document}
